\documentclass[12pt,letterpaper]{article}

% Page and typography
\usepackage[letterpaper,margin=1in]{geometry}
\usepackage{newtxtext,newtxmath}
\usepackage{setspace}
\usepackage{microtype}
\usepackage{lastpage}
\usepackage{fancyhdr}
\usepackage{tikz}

% Figures, tables, links, and citations
\usepackage{graphicx}
\usepackage{booktabs}
\usepackage{array}
\usepackage{caption}
\usepackage[hidelinks]{hyperref}
\usepackage[backend=biber,style=apa,sorting=nyt]{biblatex}
\DeclareLanguageMapping{american}{american-apa}
\addbibresource{references.bib}

% Lists and heading style
\usepackage{enumitem}
\usepackage{titlesec}

\setstretch{1.08}
\setlength{\parindent}{0pt}
\setlength{\parskip}{0.65em}
\setlength{\headheight}{15pt}
\setlist[enumerate]{leftmargin=2em,itemsep=0pt,topsep=0.2em}
\captionsetup{
  font=it,
  justification=raggedright,
  singlelinecheck=false,
  skip=4pt
}

\titleformat{\section}
  {\normalfont\bfseries\large}
  {}{0pt}{\underline}
\titleformat{\subsection}
  {\normalfont\bfseries\normalsize}
  {}{0pt}{}
\titlespacing*{\section}{0pt}{1.2em}{0.3em}
\titlespacing*{\subsection}{0pt}{0.8em}{0.2em}

% Page number and border
\pagestyle{fancy}
\fancyhf{}
\fancyfoot[R]{\thepage}
\renewcommand{\headrulewidth}{0pt}
\renewcommand{\footrulewidth}{0pt}
\AddToHook{shipout/background}{%
  \begin{tikzpicture}[remember picture,overlay]
    \draw[line width=0.4pt]
      ([xshift=0.5in,yshift=-0.5in]current page.north west)
      rectangle
      ([xshift=-0.5in,yshift=0.5in]current page.south east);
  \end{tikzpicture}%
}

\newcommand{\essaytitle}{Your Essay Title}
\newcommand{\studentname}{Your Name}
\newcommand{\studentid}{Student ID}
\newcommand{\coursename}{Course Name}
\newcommand{\submissiondate}{20 November 2026}
\newcommand{\wordcount}{0000}

\begin{document}

% The sample begins with a contents table rather than a separate title page.
\begin{center}
\begin{tabular}{>{\raggedright\arraybackslash}p{0.56\textwidth}
                >{\centering\arraybackslash}p{0.22\textwidth}}
\toprule
\textbf{Content} & \textbf{Pages} \\
\midrule
Abstract & \pageref{sec:abstract} \\
Introduction & \pageref{sec:introduction} \\
Methodology \& Questions & \pageref{sec:methodology} \\
Results \& Discussion & \pageref{sec:results}--\pageref{sec:limitations} \\
Limitations & \pageref{sec:limitations} \\
Conclusion & \pageref{sec:conclusion} \\
Reference list & \pageref{sec:references} \\
\bottomrule
\end{tabular}
\end{center}

\textit{\textbf{Words: \wordcount\ (excluding figure and table descriptions)}}

\section{Abstract}
\label{sec:abstract}
Write a concise summary of the essay. State the topic, the problem or
research focus, the approach used, the main findings, and the conclusion.
An abstract is usually one or two paragraphs.

\section{Introduction}
\label{sec:introduction}
Introduce the wider context and explain why the topic matters. Define
important terms, establish the problem, and state the purpose and scope of
the essay. Use author--date citations, for example \parencite{example2024}.

\begin{figure}[htbp]
  \centering
  % Replace this placeholder with a file in the project folder.
  % \includegraphics[width=0.82\textwidth]{figures/your-figure}
  \fbox{\parbox[c][1.6in][c]{0.78\textwidth}{\centering Insert figure here}}
  \caption{Descriptive figure title (Author, Year).}
  \label{fig:introduction-example}
\end{figure}

\section{Methodology}
\label{sec:methodology}
Explain how information was collected and evaluated. Identify whether the
essay uses a literature review, datasets, case studies, or another approach.
State the limitations of the chosen method where appropriate.

\subsection*{Questions}
\begin{enumerate}
  \item What is the first research question?
  \item What is the second research question?
  \item What is the third research question?
  \item What is the fourth research question?
\end{enumerate}

\section{Results and Discussion}
\label{sec:results}
Present and interpret the evidence. Organise this section with short
subsections that directly answer the research questions.

\subsection{First Results Theme}
Write the first results theme here. Explain what the evidence shows and
connect it to the literature.

\subsection{Second Results Theme}
Write the second results theme here. Add figures or tables when they clarify
the argument.

\begin{table}[htbp]
  \centering
  \caption{Example table title.}
  \label{tab:example}
  \begin{tabular}{ll}
    \toprule
    Category & Description \\
    \midrule
    Example 1 & Replace with your data \\
    Example 2 & Replace with your data \\
    \bottomrule
  \end{tabular}
\end{table}

\section{Limitations}
\label{sec:limitations}
Discuss the boundaries of the evidence, uncertainties in the method, data
quality, and any conclusions that should not be generalised.

\section{Conclusion}
\label{sec:conclusion}
Answer the research questions directly, summarise the main findings, and
explain the significance of the essay. Do not introduce new evidence here.

\clearpage
\section*{References}
\label{sec:references}
\addcontentsline{toc}{section}{Reference list}
\printbibliography[heading=none]

\end{document}
